\begin{table}[H]
\centering
\caption{Regime averages, Portugal ESA 2010 sample}
\label{tab:regimes}
\resizebox{\textwidth}{!}{%
\begin{tabular}{lrrrrrrr}
\toprule
Regime & Years & Int./GDP & Int. EUR m & Debt/GDP & Avg. debt rate & Nom. growth & Prim. bal. \\
\midrule
Euro convergence & 1995--1998 & 4.30 & 4284.68 & 59.95 & 6.74 & 6.74 & -0.20 \\
Early euro period & 1999--2007 & 2.82 & 4131.03 & 64.07 & 4.64 & 5.27 & -1.54 \\
Global financial crisis & 2008--2010 & 3.00 & 5349.47 & 87.67 & 3.67 & 0.85 & -5.37 \\
Sovereign-debt crisis and adjustment & 2011--2014 & 4.67 & 8033.05 & 126.47 & 3.81 & -0.91 & -1.95 \\
Low-rate and asset-purchase period & 2015--2019 & 3.70 & 7219.64 & 125.08 & 2.98 & 4.37 & 1.76 \\
Pandemic & 2020--2021 & 2.60 & 5407.45 & 129.00 & 2.05 & 0.71 & -1.70 \\
Inflation and monetary tightening & 2022--2025 & 1.98 & 5515.15 & 97.83 & 2.05 & 9.14 & 2.50 \\
\bottomrule
\end{tabular}%
}
\smallskip
\small Notes: All entries are regime means. Interest in euros is reported in million euros. Displayed values use Python's round-half-to-even convention; the equal Euro-convergence means for the average-debt rate and nominal growth are a rounding coincidence. The ten-year yield is excluded from the printed table to keep the regime table readable.
\end{table}
